\chapter{Success Criteria}

future modules. C-6 Endurance. Runtime is bounded by available battery technology within the size/mass envelope. C-7 Envelope. The robot must remain within the stated footprint, height, and mass limits to fit aisles and remain transportable. C-8 Regulatory. All trials must comply with applicable local workplace and electrical safety regulations. C-9 Schedule. Scope is intentionally limited to what is achievable within the V1 prototype timeline.

7 Success Criteria V1 is considered successful when the following measurable criteria are met under the defined test conditions. Each criterion is traceable to the requirements it validates.

ID         Success criterion (measurable)                            Traces to        Verif. SC-01      Complete a predefined route of $\geq$50 m autonomously         FR-                 T in at least 8 of 10 runs ( $\geq$80\% success ).                09,10,11 SC-02      Carry the 50 kg rated payload over the full route with    FR-04,19            T no loss of mobility, traction, or stability. SC-03      Across $\geq$20 obstacle encounters, stop before contact       FR-14,15;           T every time (zero collisions).                             SR-04

ID         Success criterion (measurable)                          Traces to        Verif. SC-04      Demonstrate controlled motion across the full           FR-03;              T 0.5--1.5 m/s range without exceeding 1.5 m/s.            SR-05 SC-05      E-stop brings the robot to a complete stop within       SR-                 T 0.5 s in every operating mode.                          01,02,12 SC-06      Operate continuously for $\geq$3 h on a single charge        FR-22               T under nominal load.

